% =====================================================================
% HPRC 2026 Poster
% A pangenome sequence search and coordinate translation service
% for the UCSC Genome Browser
%
% Parsa Eskandar, Jouni Siren, Benedict Paten
% UC Santa Cruz Genomics Institute
%
% Build: Overleaf (pdfLaTeX) or tectonic.  48in x 36in landscape.
% =====================================================================
\documentclass[final]{beamer}
\usepackage[orientation=landscape,size=custom,width=121.92,height=91.44,scale=1.0]{beamerposter}

\usepackage[T1]{fontenc}
\usepackage[sfdefault]{FiraSans}
\usepackage{amsmath}
\usepackage{microtype}
\usepackage{graphicx}
\usepackage{ragged2e}
\usepackage{tikz}
\usetikzlibrary{arrows.meta,positioning,calc,decorations.pathreplacing}
\usepackage{pgfplots}
\pgfplotsset{compat=1.17}
\usepackage[most]{tcolorbox}
\usepackage{qrcode}
\usepackage{fontawesome5}

% ---------------------------------------------------------------------
% Palette
% ---------------------------------------------------------------------
\definecolor{navy}{HTML}{003C6C}      % UCSC primary
\definecolor{navydeep}{HTML}{06294A}  % masthead
\definecolor{ink}{HTML}{18293B}       % body text
\definecolor{muted}{HTML}{5B6B7C}     % secondary text
\definecolor{gold}{HTML}{FDC700}      % UCSC gold
\definecolor{amber}{HTML}{C98A00}     % gold, dark enough for text
\definecolor{teal}{HTML}{0E7C9C}
\definecolor{coral}{HTML}{E4572E}
\definecolor{green}{HTML}{2E7D32}
\definecolor{pagebg}{HTML}{EDF1F5}
\definecolor{hairline}{HTML}{D9E1E9}
\definecolor{chipbg}{HTML}{F1F4F8}
\definecolor{skyfill}{HTML}{E7F0F7}

\setbeamercolor{background canvas}{bg=pagebg}
\setbeamertemplate{navigation symbols}{}
\setbeamersize{text margin left=0cm, text margin right=0cm}
\setbeamercolor{normal text}{fg=ink}
\setbeamercolor{item}{fg=navy}

% ---------------------------------------------------------------------
% Card system
% ---------------------------------------------------------------------
\newtcolorbox{card}{
  enhanced, colback=white, colframe=white, boxrule=0pt, arc=5mm,
  left=1.1cm, right=1.1cm, top=0.75cm, bottom=0.85cm,
  before skip=0cm, after skip=0.95cm,
  fuzzy shadow={0mm}{-2mm}{0mm}{0.55mm}{black!20}
}

% \cardhead{NN}{Title}
\newcommand{\cardhead}[2]{%
  \begin{tikzpicture}[baseline=(t.base)]
    \node[fill=gold, rounded corners=2.2mm, inner sep=0pt,
          minimum width=1.7cm, minimum height=1.7cm,
          font=\large\bfseries, text=navydeep] (n) {#1};
    \node[right=0.6cm of n.east, anchor=west, inner sep=0pt,
          font=\Large\bfseries, text=navy] (t) {#2};
  \end{tikzpicture}\par
  \vspace{0.4cm}
  {\color{hairline}\hrule height 1.5pt}
  \vspace{0.6cm}
}

% Small stat chip used in the hero strip
\newcommand{\statchip}[2]{%
  \begin{tikzpicture}
    \node[fill=chipbg, rounded corners=3mm, inner xsep=0.75cm, inner ysep=0.38cm,
          align=center] {%
      {\color{navy}\fontsize{40}{42}\selectfont\bfseries #1}\\[0.35ex]
      {\color{muted}\small #2}};
  \end{tikzpicture}}

% Inline label chip
\newcommand{\chip}[1]{\tikz[baseline=(c.base)]{\node[fill=chipbg, rounded corners=1.8mm,
  inner xsep=0.4cm, inner ysep=0.18cm, font=\small, text=ink] (c) {#1};}}

% Icon bullet for capability / future-work rows
\newcommand{\iconrow}[3]{% {icon}{bold lead}{body}
  \par\noindent
  \begin{minipage}[t]{2.1cm}
    \centering\tikz{\node[fill=skyfill, circle, minimum size=1.8cm,
      text=navy, font=\large]{#1};}
  \end{minipage}\hfill
  \begin{minipage}[t]{\dimexpr\linewidth-2.6cm\relax}
    \RaggedRight{\bfseries\color{navy}#2}\\[0.15ex]#3
  \end{minipage}\par\vspace{0.5cm}}

% Browser-chrome frame for screenshots (placeholder until PNG exists)
\newcommand{\browserframe}[3]{% {filename}{height}{caption}
  \begin{center}
  \begin{tikzpicture}
    \node[fill=chipbg, rounded corners=3mm, inner sep=0pt] (win)
      {\begin{minipage}{0.94\linewidth}
        \vspace{0.3cm}
        \hspace{0.5cm}\tikz{\fill[coral!80] (0,0) circle (0.13);
          \fill[gold] (0.5,0) circle (0.13);
          \fill[green!70] (1.0,0) circle (0.13);}%
        \hspace{0.6cm}{\footnotesize\color{muted}\texttt{\detokenize{genome.ucsc.edu}}}
        \vspace{0.25cm}\par
        \IfFileExists{figures/#1}{%
          \includegraphics[width=\linewidth,height=#2,keepaspectratio]{figures/#1}\par
        }{%
          \begin{tikzpicture}
            \node[fill=white, minimum width=\linewidth, minimum height=#2,
                  align=center, text=muted, font=\normalsize, inner sep=0.4cm]
              {\faImage\ \ add \texttt{\detokenize{figures/#1}}\\[0.4ex]
               {\footnotesize #3}};
          \end{tikzpicture}\par
        }%
      \end{minipage}};
  \end{tikzpicture}
  \end{center}}

% Chart defaults
\pgfplotsset{
  poster/.style={
    axis line style={hairline, thick},
    axis x line*=bottom, axis y line*=left,
    tick style={hairline, thick},
    ticklabel style={font=\small, color=muted},
    label style={font=\small, color=muted},
    grid style={hairline, dashed},
    every axis/.append style={line width=1.2pt},
  },
  sqlegend/.style={legend image code/.code={\fill[##1] (0cm,-0.1cm) rectangle (0.38cm,0.26cm);}},
}

\begin{document}
\begin{frame}[t]

% ====================================================================
%  MASTHEAD
% ====================================================================
\vspace*{-0.15cm}
\begin{tikzpicture}
  \fill[navydeep] (0,0) rectangle (\paperwidth,-12.5cm);
  \fill[gold]     (0,-12.5cm) rectangle (\paperwidth,-12.85cm);

  % --- emblem: haplotype ribbons through two bubbles -----------------
  \begin{scope}[shift={(3.0cm,-6.2cm)}, xscale=1.28, yscale=1.28,
                line width=9pt, line cap=round]
    \draw[white!85!navy]  (0,0)   .. controls (2.2,0)   and (2.8,1.5)  .. (4.6,1.5)
                          .. controls (6.4,1.5) and (7.0,0)   .. (9.2,0)
                          .. controls (11.4,0)  and (12.0,0)  .. (13.6,0);
    \draw[gold]           (0,0)   .. controls (2.6,0)   and (3.2,-1.5) .. (4.6,-1.5)
                          .. controls (6.0,-1.5) and (6.8,0) .. (9.2,0)
                          .. controls (10.6,0) and (11.2,1.35) .. (12.2,1.35)
                          .. controls (13.0,1.35) and (13.2,0.6) .. (13.6,0);
    \draw[teal!60!white]  (0,0) -- (13.6,0);
    \fill[white] (0,0) circle (0.3) (13.6,0) circle (0.3);
  \end{scope}

  % --- title block ---------------------------------------------------
  \node[anchor=north, align=center, text=gold, font=\normalsize\bfseries]
    at (0.5\paperwidth,-1.35cm)
    {HPRC 2026 \ \textbullet\ \ HUMAN PANGENOME REFERENCE CONSORTIUM};
  \node[anchor=north, align=center, text=white,
        font=\fontsize{78}{88}\selectfont\bfseries]
    at (0.5\paperwidth,-2.55cm)
    {A Pangenome Sequence Search and Coordinate\\
     Translation Service for the UCSC Genome Browser};
  \node[anchor=north, align=center, text=white, font=\LARGE]
    at (0.5\paperwidth,-9.35cm)
    {Parsa Eskandar \quad\textbullet\quad Jouni Sir\'en \quad\textbullet\quad Benedict Paten};
  \node[anchor=north, align=center, text=white!72!navy, font=\Large]
    at (0.5\paperwidth,-11.05cm)
    {UC Santa Cruz Genomics Institute \quad\textbullet\quad \texttt{seeskand@ucsc.edu}};

  % --- QR tile --------------------------------------------------------
  \node[anchor=east, fill=white, rounded corners=3mm, inner sep=0.5cm]
    (qr) at (\paperwidth-3.0cm,-5.4cm) {\qrcode[height=5.6cm]{https://github.com/parsaeskandar/hprc_2026_poster}};
  \node[anchor=north, text=white!72!navy, font=\small, align=center]
    at (qr.south) [yshift=-0.3cm] {poster sources \ \faGithub\\[-0.1ex]\texttt{\detokenize{parsaeskandar/hprc_2026_poster}}};
\end{tikzpicture}

% ====================================================================
%  HERO STRIP
% ====================================================================
\vspace{0.55cm}
\begin{center}
\begin{minipage}{0.965\paperwidth}
\begin{tcolorbox}[enhanced, colback=white, boxrule=0pt, arc=5mm,
    left=1.4cm, right=1.4cm, top=0.65cm, bottom=0.65cm,
    fuzzy shadow={0mm}{-2mm}{0mm}{0.55mm}{black!20}]
  \begin{minipage}[c]{0.40\linewidth}
    \RaggedRight
    {\fontsize{32}{40}\selectfont\color{navy}
     Paste a \textbf{sequence} or a \textbf{position} --- get an answer from
     \textbf{every HPRC haplotype}, live, inside the browser.}
  \end{minipage}\hfill
  \begin{minipage}[c]{0.58\linewidth}
    \centering
    \statchip{464}{haplotypes scored per query}\hspace{0.35cm}
    \statchip{$\sim$20\,ms}{gene-scale liftover}\hspace{0.35cm}
    \statchip{466\,\texttimes\,466}{assembly pairs, on demand}\hspace{0.35cm}
    \statchip{15}{lines changed in kent core}
  \end{minipage}
\end{tcolorbox}
\end{minipage}
\end{center}
\vspace{0.45cm}

% ====================================================================
%  BODY --- four columns
% ====================================================================
\begin{columns}[t]

% ====================== COLUMN 1 =====================================
\begin{column}{0.2335\paperwidth}

\begin{card}
\cardhead{01}{The graph nobody could ask}
\RaggedRight
The HPRC v2.0 minigraph--Cactus graph holds \textbf{233 samples / 464
haplotypes}, yet there has been no interactive way to query it. Liftover
needs a chain file per assembly pair --- and almost none exist:
\begin{center}
\begin{tikzpicture}[scale=0.92]
  \foreach \x in {0,...,11}
    \foreach \y in {0,...,7}
      {\fill[hairline] (\x*0.62,\y*0.62) rectangle +(0.5,0.5);}
  \fill[gold] (4*0.62,3*0.62) rectangle +(0.5,0.5);
  \node[anchor=north, font=\small, text=muted, align=center]
    at (3.4,-0.45) {\textbf{chain files today}\\ 56 of $\sim$213{,}000 pairs};
  \node[font=\LARGE, text=navy] at (8.6,2.3) {$\Longrightarrow$};
  \begin{scope}[shift={(10.2,0)}]
    \foreach \x in {0,...,11}
      \foreach \y in {0,...,7}
        {\fill[navy!80] (\x*0.62,\y*0.62) rectangle +(0.5,0.5);}
    \node[anchor=north, font=\small, text=muted, align=center]
      at (3.4,-0.45) {\textbf{through the graph}\\ every pair, on demand};
  \end{scope}
\end{tikzpicture}
\end{center}
\vspace{0.15cm}
A researcher with a \textbf{sequence} can't ask which haplotypes carry it;
one with a \textbf{position} can't move it between assemblies. We built the
service that answers both, and wired it into the UCSC Genome Browser.
\end{card}

\begin{card}
\cardhead{02}{What users get}
\iconrow{\faSearch}{Pangenome sequence search}{A BLAT-like page: a pasted
sequence is mapped to the whole graph with \texttt{vg giraffe}; results rank
all 464 haplotypes by identity and link to its position on any of them,
drawn as a track with base-level mismatch coloring.}
\iconrow{\faRandom}{Any-to-any coordinate translation}{Liftover between
any two of 466 assemblies at interactive speed, with a scored ``which
assemblies contain this region'' destination picker.}
\iconrow{\faLayerGroup}{Annotation lifting}{A QuickLift chain generated
\emph{at request time}: users land on an HPRC assembly with CHM13's genes,
RefSeq, CenSat and T2T~Encode at translated positions, plus an
\textbf{Alignment Differences} track.}
\end{card}

\begin{card}
\cardhead{03}{Architecture}
\begin{center}
\begin{tikzpicture}[
  layer/.style={draw=hairline, very thick, fill=skyfill, rounded corners=3mm,
                align=center, text width=21.0cm, inner ysep=0.34cm},
  arr/.style={-{Stealth[length=5mm]}, line width=2.2pt, navy!60}
]
  \node[layer] (gb) {{\bfseries\color{navy}UCSC Genome Browser}\\[0.2ex]
    {\small \texttt{hgPangenome} CGI + JS \ \textbullet\ \ mapping page \ \textbullet\ \ convert page \ \textbullet\ \ \texttt{hgConvert}}};
  \node[layer, below=0.6cm of gb] (mw) {{\bfseries\color{navy}Python middleware} {\small(JSON API, token auth)}\\[0.25ex]
    {\small\texttt{/map} \ \ \texttt{/surject} \ \ \texttt{/liftover} \ \ \texttt{/liftover/targets}}};
  \node[layer, below=0.6cm of mw.south west, anchor=north west, text width=10.0cm] (vg)
    {{\bfseries\color{navy}vg giraffe-server}\\[0.15ex]{\small persistent mapper;\\ indexes loaded once}};
  \node[layer, below=0.6cm of mw.south east, anchor=north east, text width=10.0cm] (cpp)
    {{\bfseries\color{navy}pangenome-index}\\[0.15ex]{\small C++ translation engine}};
  \node[below=0.45cm of cpp.south, anchor=north, text width=10.0cm, align=center] (idx)
    {\small\color{muted}GBZ \ 148M nodes \ \textbullet\ \ r-index + tag array\\ FastLocate $\sim$11\,GB \ \textbullet\ \ Tables 1+2 \ 1.78\,GB};
  \draw[arr] (gb) -- (mw);
  \draw[arr] (mw.south -| vg) -- (vg);
  \draw[arr] (mw.south -| cpp) -- (cpp);
  \draw[arr, hairline] (cpp) -- (idx);
\end{tikzpicture}
\end{center}
\vspace{0.1cm}
\RaggedRight
Core browser changes: \textbf{3 files, 15 lines}. Everything else is a new,
self-contained CGI --- dark unless configured.
\end{card}

\end{column}

% ====================== COLUMN 2 =====================================
\begin{column}{0.2335\paperwidth}

\begin{card}
\cardhead{04}{Translation through the graph}
\begin{center}
\begin{tikzpicture}[line cap=round, scale=0.97]
  % ---- source bar ---------------------------------------------------
  \node[anchor=west, font=\small\bfseries, text=navy] at (0.1,12.35) {source \ \ CHM13\#0 \ chr10};
  \fill[navy!12, rounded corners=2mm] (0,11.0) rectangle (22,11.85);
  \fill[gold, rounded corners=1mm] (7,11.0) rectangle (13.5,11.85);
  \node[font=\footnotesize, text=navydeep] at (10.25,11.42) {query interval};
  \node[fill=gold, circle, minimum size=1.0cm, font=\small\bfseries, text=navydeep] at (5.7,11.42) {1};

  % ---- probes annotation (top) ---------------------------------------
  \node[fill=gold, circle, minimum size=1.0cm, font=\small\bfseries, text=navydeep] at (0.7,9.0) {2};
  \node[anchor=west, font=\footnotesize, text=amber] at (1.5,9.0)
    {probes: 1 / 2\,kb --- each returns \emph{every} haplotype on the node};

  % ---- graph row ------------------------------------------------------
  \coordinate (n1) at (1.5,5.2);  \coordinate (n2) at (4.0,5.2);
  \coordinate (n3) at (6.5,5.2);  \coordinate (bU) at (9.0,6.3);
  \coordinate (bD) at (9.0,4.1);  \coordinate (n4) at (11.5,5.2);
  \coordinate (n5) at (14.0,5.2); \coordinate (n6) at (16.5,5.2);
  \coordinate (n7) at (19.0,5.2); \coordinate (n8) at (21.2,5.2);
  \coordinate (det) at (17.75,4.45);
  \draw[hairline, line width=2.4pt]
    (n1)--(n2)--(n3) (n4)--(n5)--(n6) (n7)--(n8)
    (n3) to[bend left=20] (bU) (bU) to[bend left=20] (n4)
    (n3) to[bend right=20] (bD) (bD) to[bend right=20] (n4)
    (n6) to[bend left=28] (n7)
    (n6) -- (det) -- (n7);
  \draw[teal, line width=4.2pt, opacity=0.7]
    (0.5,5.32) -- (n1) -- (n2) -- (n3) to[bend left=20] (bU) to[bend left=20] (n4)
    -- (n5) -- (n6) to[bend left=28] (n7) -- (n8) -- (22,5.32);
  \draw[coral, line width=4.2pt, opacity=0.7]
    (0.5,5.08) -- (n1) -- (n2) -- (n3) to[bend right=20] (bD) to[bend right=20] (n4)
    -- (n5) -- (n6) -- (det) -- (n7) -- (n8) -- (22,5.08);
  \foreach \p in {n1,n2,n3,bU,bD,n4,n5,n6,n7,n8}
    {\fill[white] (\p) circle (0.40); \draw[navy!55, line width=2.2pt] (\p) circle (0.40);}
  \fill[navy!30] (det) circle (0.26);
  % probe arrows
  \foreach \p in {n2,n5,n7}
    {\draw[amber, line width=2.4pt, -{Stealth[length=4mm]}] ($(\p)+(0,2.5)$) -- ($(\p)+(0,0.62)$);}
  % anchors
  \draw[gold, line width=4pt] (n3) circle (0.58);
  \draw[gold, line width=4pt] (n6) circle (0.58);
  \node[fill=gold, circle, minimum size=1.0cm, font=\small\bfseries, text=navydeep] at (0.7,2.6) {3};
  \node[anchor=west, font=\footnotesize, text=ink] at (1.5,2.6)
    {unique anchors {\color{muted}(1$\times$ on both)}};
  % trace brace (below)
  \draw[decorate, decoration={brace, mirror, amplitude=8pt}, navy, line width=2pt]
    (11.6,3.6) -- (15.9,3.6);
  \node[fill=gold, circle, minimum size=1.0cm, font=\small\bfseries, text=navydeep] at (12.2,2.6) {4};
  \node[anchor=west, font=\footnotesize, text=navy] at (13.0,2.6)
    {LF-walk trace $\to$ offsets};

  % ---- target bar ------------------------------------------------------
  \fill[teal!12, rounded corners=2mm] (0,0.0) rectangle (22,0.85);
  \fill[gold, rounded corners=1mm] (9,0.0) rectangle (15.5,0.85);
  \node[font=\footnotesize, text=navydeep] at (12.25,0.42) {translated interval};
  \node[anchor=west, font=\small\bfseries, text=teal] at (0.1,-0.55) {target \ \ HG00408\#1};
  % projections
  \draw[dashed, amber, line width=1.8pt] (7,11.0) -- (6.5,5.9);
  \draw[dashed, amber, line width=1.8pt] (13.5,11.0) -- (16.5,5.9);
  \draw[dashed, amber, line width=1.8pt] (6.5,4.5) -- (9,0.85);
  \draw[dashed, amber, line width=1.8pt] (16.5,4.5) -- (15.5,0.85);
\end{tikzpicture}
\end{center}
\vspace{0.35cm}
\RaggedRight
{\bfseries\color{navy}1} \ Table 1 turns \texttt{contig:start-end} into a
graph path interval. \ {\bfseries\color{navy}2} \ Probing discovers target
candidates --- or the \textbf{routing Table 2} skips it with a direct
lookup. \ {\bfseries\color{navy}3} \ Nodes visited exactly once by both
sides are guaranteed orthologous; a colinearity gate drops implausible
pairs. \ {\bfseries\color{navy}4} \ The trace folds into
\textbf{chain-semantics blocks}: broken by indels, flips, contig changes ---
never by substitutions.
\end{card}

\begin{card}
\cardhead{05}{A routing table that fits}
\RaggedRight
An all-pairs interval table over the frequency-filtered graph (224M path
fragments, $2{,}367\times$ fragmentation) would cost \textbf{17.2\,TB}.
Building on the \emph{non-filtered} graph (94{,}554 paths) and storing
\emph{routing only} --- which targets a source range touches, never their
coordinates:
\begin{center}
\begin{tikzpicture}
\begin{axis}[poster, width=0.9\linewidth, height=6.2cm,
  xbar, xmode=log, xmin=0.5, xmax=400000,
  symbolic y coords={routing --- ours, star via references, all-pairs filtered},
  ytick=data, bar width=0.95cm,
  enlarge y limits=0.3,
  xlabel={index size (GB, log scale)},
  nodes near coords style={font=\small\bfseries, color=ink},
  point meta=explicit symbolic, nodes near coords,
  ]
  \addplot[fill=navy!80, draw=none] coordinates {
    (17610,all-pairs filtered) [17.2 TB]
    (74,star via references) [74 GB]
    (1.78,routing --- ours) [1.78 GB]
  };
\end{axis}
\end{tikzpicture}
\end{center}
Dropping stored target intervals also removed \emph{every} failure mode of
interval pairing:\\[0.3cm]
\chip{\textbf{0} abandoned runs \ {\color{muted}(was 5.0M)}}\
\chip{\textbf{0} splits \ {\color{muted}(was 35.7M)}}\
\chip{build: \textbf{35 min}}
\end{card}

\begin{card}
\cardhead{06}{Scoring all 464 at once}
\RaggedRight
One pass over a mapped read's nodes scores \textbf{every haplotype}
simultaneously --- each probe already returns all of them:
\begin{center}
\begin{tikzpicture}[line cap=round, scale=0.97]
  \node[anchor=west, font=\small, text=muted] at (0,3.55) {aligned read bases};
  \fill[navy!25, rounded corners=1.5mm] (0,2.65) rectangle (21,3.3);
  \node[anchor=west, font=\small, text=muted] at (0,2.1) {on nodes haplotype $H$ visits};
  \fill[teal!75, rounded corners=1.5mm] (0,1.1) rectangle (17.6,1.75);
  \fill[coral] (4.1,1.1) rectangle (4.55,1.75);
  \fill[coral] (11.2,1.1) rectangle (11.8,1.75);
  \node[anchor=west, font=\footnotesize, text=coral] at (17.9,1.42) {mismatches};
  \draw[decorate, decoration={brace, mirror, amplitude=9pt}, navy, line width=2pt]
    (0,0.75) -- (17.6,0.75);
  \node[font=\small\bfseries, text=navy] at (8.8,-0.2)
    {coverage$(H)$ = 84\% \quad\textbullet\quad identity$(H)$ = 79\%};
\end{tikzpicture}
\end{center}
\vspace{0.2cm}
identity $\le$ coverage always; the gap is the mismatch burden. Together
they separate cases one number conflates:\\[0.3cm]
\chip{100 / 99 \ \ present, near-exact}\
\chip{100 / 60 \ \ present, diverged}\
\chip{50 / 50 \ \ half missing}
\end{card}

\end{column}

% ====================== COLUMN 3 =====================================
\begin{column}{0.2335\paperwidth}

\begin{card}
\cardhead{07}{Sequence search, in the browser}
\browserframe{mapping_page.png}{8.6cm}{Pangenome Mapping page: identity-ranked assemblies, position-on-haplotype links}
\RaggedRight
A 300\,bp query scores all 464 haplotypes in seconds. Picking a haplotype
\textbf{re-surjects the cached alignment} --- no remapping --- and the
browser opens there with the sequence as a native track:
\browserframe{pangenome_seq_track.png}{6.2cm}{``Pangenome Seq'' track, BLAT-style base-level mismatch coloring}
\end{card}

\begin{card}
\cardhead{08}{Landing with annotations}
\browserframe{lifted_annotations.png}{8.6cm}{HPRC assembly with lifted CHM13 tracks + Alignment Differences}
\RaggedRight
Converting a position lands on the target with \textbf{CHM13's annotation
tracks lifted} through an on-the-fly chain, plus an \textbf{Alignment
Differences} track --- verified base-by-base: across 22\,kb of ABO it shows
exactly the 2 real C$\to$T mismatches between hs1 and HG00597\#1. Pangenome
assemblies also appear in the stock \texttt{hgConvert} dropdown.
\end{card}

\begin{card}
\cardhead{09}{Chains that keep up with panning}
\begin{center}
\begin{tikzpicture}[line cap=round, scale=0.97]
  \draw[hairline, line width=3pt] (0,0) -- (22,0);
  \fill[gold] (1.2,-0.2) rectangle (4.4,0.2);
  \fill[teal!70] (8.2,-0.2) rectangle (17.8,0.2);
  \foreach \x/\l in {1.2/click, 4.4/page loads, 8.2/beacon, 17.8/done}
    {\draw[ink, line width=2pt] (\x,-0.45) -- (\x,0.45);
     \node[font=\footnotesize, text=muted, anchor=north] at (\x,-0.6) {\l};}
  \node[font=\small\bfseries, text=amber, anchor=south] at (2.8,0.55)
    {phase 1 \ 0.7\,s};
  \node[font=\footnotesize, text=amber, anchor=south] at (2.8,1.55)
    {pad to 5\texttimes\ span};
  \node[font=\small\bfseries, text=teal, anchor=south] at (13.0,0.55)
    {phase 2 \ 11\,s in background};
  \node[font=\footnotesize, text=teal, anchor=south] at (13.0,1.55)
    {222-byte \texttt{sendBeacon} $\to$ rebuild at \textpm 5\,Mb};
  \node[font=\footnotesize, text=muted, anchor=west] at (18.4,0.85)
    {next pan: free};
\end{tikzpicture}
\end{center}
\vspace{0.4cm}
\RaggedRight
Blocks become bigChain{+}bigLink (JSON $\to$ BED $\to$
\texttt{bedToBigBed}); QuickLift reopens the file each request, so the
widened chain appears on the next pan --- \emph{no reload}. Lifted items
visible after panning away:
\begin{center}
\begin{tikzpicture}
\begin{axis}[poster, width=0.92\linewidth, height=7.2cm,
  ybar, ymode=log, ymin=0.5, ymax=600,
  symbolic x coords={in window, 400 kb, 2 Mb, 4 Mb},
  xtick=data, bar width=1.35cm,
  enlarge x limits=0.16,
  ylabel={lifted gene items (log)},
  legend style={font=\small, draw=none, fill=none, at={(0.5,1.02)}, anchor=south,
    /tikz/every even column/.append style={column sep=0.9cm}},
  legend columns=2, legend cell align=left, sqlegend,
  point meta=explicit symbolic,
  nodes near coords, every node near coord/.append style={font=\small\bfseries, color=ink},
  ]
  \addplot[fill=gold, draw=none] coordinates
    {(in window,3) [3] (400 kb,129) [129] (2 Mb,1) [1] (4 Mb,1) [1]};
  \addplot[fill=navy!80, draw=none] coordinates
    {(in window,3) [3] (400 kb,129) [129] (2 Mb,5) [5] (4 Mb,31) [31]};
  \legend{phase 1 (1.1 Mb chain), phase 2 (10 Mb chain)}
\end{axis}
\end{tikzpicture}
\end{center}
\end{card}

\end{column}

% ====================== COLUMN 4 =====================================
\begin{column}{0.2335\paperwidth}

\begin{card}
\cardhead{10}{Performance}
\RaggedRight
\textbf{Translation latency} (warm; completeness in gray):
\begin{center}
\begin{tikzpicture}
\begin{axis}[poster, width=0.9\linewidth, height=8.4cm,
  xbar, xmode=log, xmin=8, xmax=300000,
  symbolic y coords={1 Mb, 100 kb, 10 kb, 1 kb, 100 bp},
  ytick=data, bar width=0.95cm,
  enlarge y limits=0.14,
  xlabel={median latency (ms, log scale)},
  point meta=explicit symbolic,
  nodes near coords, nodes near coords style={font=\small\bfseries, color=ink},
  ]
  \addplot[fill=navy!80, draw=none] coordinates {
    (22,100 bp)  [22 ms \ {\color{muted}99.0\%}]
    (21,1 kb)    [21 ms \ {\color{muted}99.8\%}]
    (115,10 kb)  [115 ms \ {\color{muted}99.6\%}]
    (599,100 kb) [0.6 s \ {\color{muted}97.2\%}]
    (4300,1 Mb)  [4.3 s \ {\color{muted}93.1\%}]
  };
\end{axis}
\end{tikzpicture}
\end{center}
Exhaustive candidate discovery (was stop-at-first-hit) took 1\,Mb
completeness \textcolor{coral}{\textbf{2.9\%}} $\to$
\textcolor{green}{\textbf{93.1\%}} --- and made it $2\times$ faster.
\vspace{0.45cm}

\textbf{Concurrent queries} (3.4\,kb, GIL released in C++):
\begin{center}
\begin{tikzpicture}
\begin{axis}[poster, width=0.9\linewidth, height=7.8cm,
  xmode=log, log basis x=2, xtick={1,2,4,8,16,32},
  xticklabels={1,2,4,8,16,32},
  ymin=0, ymax=80,
  ytick={0,20,40,60,80}, yticklabels={0,20,40,60,80},
  xlabel={threads}, ylabel={queries / s},
  grid=major,
  ]
  \addplot[navy, line width=3pt, mark=*, mark size=3pt,
           mark options={fill=navy}] coordinates
    {(1,6.3) (2,13.2) (4,23.6) (8,60.3) (16,59.4) (32,65.3)};
  \draw[gold, line width=3pt] (axis cs:8,60.3) circle (0.32cm);
  \node[font=\small\bfseries, color=amber, anchor=south east]
    at (axis cs:8,62) {9.6$\times$ at 8 threads};
\end{axis}
\end{tikzpicture}
\end{center}
Scored destination picker: 100\,kb region in \textbf{1.28\,s} vs 29.5\,s
exact ($23\times$). Chain generation: linear, $\sim$1.14\,s/Mb.
\end{card}

\begin{card}
\cardhead{11}{What's next}
\iconrow{\faBolt}{Bulk offline chain generator}{Direct path intersection
instead of r-index probing: est.\ $\sim$300$\times$ faster ($\sim$1.6\,h per
target genome-wide vs 19 days via the API) --- precomputed whole-chromosome
chains, full QuickLift parity.}
\iconrow{\faDatabase}{Persistent chain cache}{Widened chains currently die
with the one-hour session; a cache makes popular regions free after the
first visit.}
\iconrow{\faPuzzlePiece}{Segdup-dense intervals}{Wide queries lose positions
($\sim$46\% at 1q21.1); exact path intersection removes the probe spacing
that misses them.}
\iconrow{\faRocket}{Public release}{Core-change sign-off, per-install API
tokens, service-side concurrency deployment.}
\end{card}

\begin{card}
\cardhead{12}{Acknowledgements}
\RaggedRight\small
Built on \texttt{vg}, GBWT/GBZ and the tag-array index (Sir\'en et al.), the
HPRC v2.0 minigraph--Cactus graph, and the UCSC Genome Browser (kent).
\texttt{vg giraffe-server} is under review as vgteam/vg PR \#4879. We thank
the UCSC Genome Browser team for design feedback, and the Human Pangenome
Reference Consortium.
\end{card}

\end{column}

\end{columns}
\end{frame}
\end{document}
